This section interprets the results for deployment, covering manufacturing
variability, the saturation of the standard one-step benchmark, the chemistry
alignment principle, the complementary roles of the early-life fingerprint and
fine tuning, pretraining-run variance and reproducibility, and the model's
computational footprint, before setting out the study's limitations.

\subsection{Manufacturing Variability as the Primary Bottleneck}

The central lesson of the SIT data is that prediction error lives in the cell,
not the algorithm: within-cell trajectories are highly predictable, yet
nominally identical cells diverge so widely (Section~\ref{sec:variability}) that
cell identity, not the choice of method, sets the error floor
(Section~\ref{sec:baseline}). The field's effort, dominated by architectural
innovation, is thus allocated against where the variance actually lies. For a
BMS engineer the practical consequence is an asymmetric risk: same-cell models
underestimate catastrophic degraders, while naive cross-cell pooling degrades
ordinary robust cells; the fingerprint encoder is our response, conditioning
each prediction on the cell's own early-life behaviour before any fine tuning
begins. The same variability also explains why benchmarks built on hand-picked,
smoothly degrading cells are systematically optimistic, and why a dataset that
spans catastrophic to exceptionally robust cells is a more honest proving
ground.

\subsection{Benchmark Saturation: A Floor Check, Not a Leaderboard}

The most consequential methodological finding is that rolling one-step-ahead
prediction is saturated on these smooth trajectories: a trivial 20-cycle linear
autoregressor matches or beats every deep model, and inter-method gaps fall
within the run-to-run variance of training itself (Section~\ref{sec:baseline}).
The obvious harder alternative, recursive multi-step rollout, discriminates no
better here, since feeding predictions back accumulates error until every model
is unusable. One-step RMSE is therefore a floor check, not a leaderboard. We
accordingly recommend that studies on smooth-degradation data (i)~always include
a trivial persistence or linear baseline, (ii)~report multi-run variance under
deterministic seeding, and (iii)~rank methods on the capabilities the
application needs, early-life prediction, calibrated uncertainty, and
cross-chemistry robustness, rather than on fourth-decimal RMSE. The same
discipline applies across operating conditions: apparent temperature-transfer
effects here reflect test-set difficulty rather than easy transfer, and must be
controlled for (Section~\ref{sec:brittleness}).

\subsection{The Chemistry Alignment Principle}

Uncertainty calibration transferred only when the pretraining and deployment
chemistries matched (Section~\ref{sec:cross_dataset}). The mechanism is simple:
the model's learned sense of how uncertain to be is tuned to LFP degradation
dynamics, so its intervals are appropriately scaled on LFP cells but
overconfident on the faster-fading LCO cells. The practical rule is to pretrain
on chemistry-matched data. Because even matched intervals under-cover, we regard
split-conformal correction (or temperature scaling~\cite{guo2017calibration})
as mandatory before prediction intervals are trusted in a BMS, whether or not
the chemistry is matched.

\subsection{Fingerprint and Fine Tuning: Distinct Roles, Not Rivals}

The ablation and its shuffled-fingerprint control together show that the
fingerprint's role is to anchor the LSTM in the pretrained population's regime
rather than to identify the individual cell, a weaker mechanism than we first
hypothesised but a more robust one, since it needs only a representative
pretraining population, not 30 cycles rich enough to tell cells apart. The
cross-dataset experiment then separates the two levers by role: fine tuning is
the \emph{dependable} mechanism, positive in every transfer scenario including
each cross-chemistry pair, whereas zero-shot fingerprint conditioning is the
\emph{early-life} mechanism, the only option before any history exists and most
reliable exactly when chemistry and corpus breadth match, the setting a
manufacturer qualifying a product line can guarantee. For a BMS this means a
chemistry-matched model is useful from the first cycles at no adaptation cost,
with fine tuning taking over as history accumulates and becoming indispensable
once chemistries differ.

\subsection{Pretraining-Run Variance and Reproducibility}

A surprising hazard is how strongly transfer quality depends on the pretraining
draw, and that no source-domain metric we tested identified the poorly
transferring runs (Section~\ref{sec:selection}); one-step holdout RMSE was in
fact inverted, rewarding collapsed persistence-like solutions. Any transfer
pipeline whose pretraining objective admits a trivial shortcut may share this
silent failure mode, which single-run studies cannot detect; our
calibration-cell protocol is a deployable safeguard, and we recommend that
battery-prognostics transfer studies report multi-run variance. The companion
discipline is determinism: every number in this paper reproduces byte-for-byte
from a fresh clone (Section~\ref{sec:repro}), a practice we adopted after finding
that library-default seeding alone could alter experimental conclusions.

\subsection{Computational Footprint}

A recurring justification for this work is BMS deployment, which constrains model
size and inference cost, so we report both.
The hybrid is lightweight: 21\,089 parameters, approximately 82\,KB in float32
and 41\,KB in float16, small enough to reside in the memory of an embedded BMS
microcontroller.
A point SOH estimate is a single forward pass; the 50-pass MC-Dropout
uncertainty estimate, when the passes are evaluated as one batched forward
(replicating the input across the batch dimension), completes in about 4\,ms on
both a desktop CPU and the RTX 2080 Ti.
Because SOH is updated once per charge-discharge cycle (tens of minutes to
hours), this leaves a latency margin of several orders of magnitude; even an
unoptimised implementation that runs the 50 passes sequentially
($\approx$1\,s) would be comfortably within budget.
These figures were measured on a desktop CPU rather than production embedded
silicon, so absolute timings would differ on a microcontroller; nonetheless the
41\,KB footprint and few-millisecond batched inference indicate the model fits
well within the resource envelope of contemporary BMS hardware.

\subsection{Limitations}

The findings reported above should be read alongside the boundaries of the
present study. The limitations listed below stem mainly from dataset scale and
the scope of the experimental design rather than from the hybrid framework
itself, and each points to a concrete opportunity for future work: larger cell
cohorts would sharpen the statistical claims, shorter fingerprint windows and
format-matched pretraining data would broaden applicability, and post hoc
calibration would close the residual coverage gap. We state them explicitly so
that practitioners can judge how far the present conclusions transfer to their
own deployment setting.

\begin{enumerate}
  \item \textit{G1 sample size}: With ten G1 cells, LOOCV results have limited
        statistical power. The observed two-sided pattern (pooling rescues
        catastrophic degraders and harms robust cells) is directionally
        consistent, but quantitative conclusions should be interpreted with
        this sample size in mind.

  \item \textit{Fingerprint window}: The 30-cycle fingerprint requires at least
        30 discharge cycles before cell-specific predictions are available.
        For applications requiring earlier prediction (e.g., first 10 cycles),
        shorter window studies would be needed.

  \item \textit{Residual calibration gap}: Even after split-conformal
        correction, coverage reaches 93.9\%, not the nominal 95\%, and two
        extreme cells (001-8, 003-5) remain substantially under-covered.
        Conformal guarantees assume exchangeability between calibration and
        test cells, which manufacturing outliers violate by definition;
        detecting such cells online (e.g.\ from interval-width anomalies)
        remains open.

  \item \textit{Format mismatch}: The MIT pretraining cells are 18650 cylindrical
        (1.1\,Ah); the SIT cells are 50\,Ah prismatic.
        Despite this format difference, the LFP chemistry match appears sufficient
        for meaningful transfer.
        Format-matched pretraining data would likely reduce the remaining gap further.

  \item \textit{LISHEN source corpus size}: The LISHEN dataset contributes only
        three cells to the cross-dataset evaluation (scenarios M and N).
        The consistently weak gains in these scenarios ($-$24.6\% to
        $+$14.6\% for ZS+FP) are attributed to insufficient source diversity
        rather than a fundamental limitation of the fingerprint mechanism.
        Future work should re-evaluate with a larger LISHEN corpus.

  \item \textit{Calibration-cell requirement}: The checkpoint-selection
        protocol requires three target-domain cells cycled to (or near) end of
        life. This is realistic for a manufacturer qualifying a product, but
        not for retrofitting arbitrary third-party packs; in that setting a
        chemistry-matched public dataset could serve as a weaker substitute,
        which we have not evaluated.

  \item \textit{Comparison scope}: Beyond the four LSTM-family predictors
        (Simple LSTM, PA-LSTM, GA-LSTM, DE-LSTM), we benchmarked two modern
        non-recurrent architectures (TCN, Transformer) and two trivial
        baselines under the identical protocol
        (Table~\ref{tab:method_comparison}).
        The modern baselines were run under fixed configurations rather than
        the per-cell hyperparameter search afforded to GA/PA/DE-LSTM; given
        the saturation of the one-step metric we do not expect tuning to
        change the conclusions, but other model classes (e.g.\
        Gaussian-process regression) remain for future work.

  \item \textit{Evaluation protocol}: The saturation finding implies that
        one-step RMSE, including our own accuracy numbers, has limited power
        to rank methods on these data.
        Harder protocols (fixed early-life budgets, multi-step horizons with
        scheduled sampling, knee-onset prediction) are the right arena for
        future accuracy claims; we did not re-benchmark all methods under
        such protocols here.
\end{enumerate}
